% !TEX root = ../main.tex

\chapter{La integral de Itô}
\label{chap:weiner}

\section{Procesos de estocásticos}
\label{sec:weiner:estocasticos}

\begin{definicion}[Proceso estocástico]
    \label{def:proceso_estocastico}
    Un \termino[proceso!estocástico]{proceso estocástico} es una función $X\colon[0,T]\times\Omega\to\R$ que cumple los siguientes requisitos:
    \begin{itemize}
        \item $\Omega$ es un espacio de probabilidad, luego tiene asociada una $\sigma$-álgebra $\mathcal{F}$ y una medida de probabilidad $\Pro$.
        \item $X$ es una función medible, considerando $[0,T]\times\Omega$ como el espacio de medida equipado con la $\sigma$-álgebra producto $\mathcal{B}([0,T])\otimes\mathcal{F}$ y la medida producto $\lambda\otimes\Pro$, donde $\mathcal{B}([0,T])$ son los borelianos de $[0,T]$ y $\lambda$ la medida de Lebesgue sobre $[0,T]$.
    \end{itemize}
\end{definicion}

\begin{observacion}
    Si $X\colon[0,T]\times\Omega\to\R$ es un proceso estocástico:
    \begin{itemize}
        \item Fijado un $t\in[0,T]$, la función $X_t\colon\Omega\to\R$ definida por $\omega\mapsto X(t,\omega)$ es una variable aleatoria, es decir, una función medible respecto de la $\sigma$-álgebra $\mathcal{F}$.
        \item Fijado un $\omega\in\Omega$, a la función $X_\omega\colon[0,T]\to\R$ definida por $t\mapsto X(t,\omega)$ se la llama \termino[trayectoria!de un proceso estocástico]{trayectoria} de $X$ asociada a $\omega$.
    \end{itemize}
\end{observacion}

\begin{definicion}[Filtración]
    \label{def:filtracion}
    Una \termino{filtración} es una familia creciente $\{\mathcal{F}_i\}_{i\in I}$ de sub-$\sigma$-álgebras de una $\sigma$-álgebra $\mathcal{F}$ dada, donde $I\subset\R$.

    Esto es, dados $t_1$, $t_2\in I$ tales que $t_1\leq t_2$, se cumple que $\mathcal{F}_{t_1}\subset\mathcal{F}_{t_2}\subset\mathcal{F}$.
\end{definicion}

\begin{definicion}[Proceso adaptado]
    \label{def:proceso_adaptado}
    Un proceso estocástico se dice \termino[proceso!estocástico!adaptado]{adaptado} a la filtración $\{\mathcal{F}_t\}_{t\in[0,T]}$ si para cada $t\in[0,T]$, la variable aleatoria $X_t$ es medible respecto de $\mathcal{F}_t\subset\mathcal{F}$.
\end{definicion}
\subsection{Procesos progresivamente medibles}
\begin{definicion}[Medibilidad progresiva]
    Decimos que un proceso estocástico $X\colon[0,T]\times\Omega\to\R$ adaptado a una filtración $\{\mathcal{F}_t\}_{t\in[0,T]}$ es \termino[proceso!estocástico!progresivamente medible]{progresivamente medible} si para todo $t\in[0,T]$ se cumple que el proceso restringido $X\colon[0,t]\times\Omega\to\R$ es medible respecto de $\mathcal{B}([0,t])\otimes\mathcal{F}_t$.
\end{definicion}

\begin{proposicion}[condición suficiente para la medibilidad progresiva]
    Todo proceso estocástico $X\colon[0,T]\times\Omega\to\R$ adaptado a una filtración $\{\mathcal{F}_t\}_{t\in[0,T]}$ con trayectorias continuas por la derecha (o por la izquierda) es progresivamente medible.
\end{proposicion}
\begin{proof}
    Dado un $t\in[0,T]$ arbitrario, vamos a particionar $[0,t]$ en $m$ intervalos de la misma longitud, para lo cual definimos los puntos $t_k^{(m)}:=\frac{kt}{m}$ con $k\in\{0,\dotsc,m\}$.

    De ahora en adelante, asumiremos que las trayectorias de $X$ son continuas por la derecha. Con mínimas modificaciones, se puede obtener una demostración para el caso en el que las trayectorias son continuas por la izquierda.

    Definimos el proceso escalonado
    \begin{equation}
        X^{(m)}(s,\omega):=X(0,\omega)1_{\{0\}}(s)+\sum_{k=1}^{m}X(t_{k}^{(m)},\omega)1_{(t_{k-1}, t_{k}]}(s)
    \end{equation}
    Sobre este proceso cabe señalar algunas cosas:
    \begin{itemize}
        \item Las funciones indicadoras son medibles respecto de $\mathcal{B}([0,t])$.
        \item Las variables aleatorias $X_{t_{k}}$ son medibles respecto de $\mathcal{F}_{t_k}\subset\mathcal{F}_t$.
        \item El producto de una función medible respecto de $\mathcal{B}([0,t])$ y de una función medible respecto de $\mathcal{F}_t$ es una función medible respecto de $\mathcal{B}([0,t])\otimes\mathcal{F}_t$.
    \end{itemize}
    En conclusión, $X^{(m)}$ es medible respecto de $\mathcal{B}([0,t])\otimes\mathcal{F}_t$.

    Fijado $(s,\omega)\in[0,t]\times\Omega$, existe un único $t_{k}^{(m)}$, al que llamaremos $t_{k}^{(m)}(s)$ se cumple que $s\in(t_{k-1}^{(m)}(s), t_{k}^{(m)}(s)]$, y por tanto $X^{(m)}(s,\omega)=X(t_{k}^{(m)}(s),\omega)$.

    Además, sabemos que $0<t_{k}^{(m)}(s)-s\leq \frac{1}{m}$, luego, $\lim_{m\to\infty}t_{k}^{(m)}(s)=s$.

    Usando la continuidad por la derecha de las trayectorias de $X$ tenemos que
    \begin{equation}
        \lim_{m\to\infty}X^{(m)}(s,\omega)=\lim_{m\to\infty}X\left(t_{k}^{(m)}(s),\omega\right)=X\left(\lim_{m\to\infty}t_{k}^{(m)}(s),\omega\right)=X(s,\omega)\text{. }
    \end{equation}

    Recapitulando, tenemos una sucesión $\{X^{(m)}\}_{m=1}^{\infty}$ de funciones medibles respecto de $\mathcal{B}([0,t])\otimes\mathcal{F}_t$, luego su límite, que que es el proceso restringido $X\colon[0,t]\times\Omega\to\R$ también es medible respecto de $\mathcal{B}([0,t])\otimes\mathcal{F}_t$, que es lo que queríamos demostrar.
\end{proof}
\begin{teorema}[medibilidad progresiva en $L^2$]
    Todo proceso estocástico $X\colon[0,T]\times\Omega\to\R$ en $L^2([0,T]\times\Omega)$ y adaptado a una filtración $\{\mathcal{F}_t\}_{t\in[0,T]}$ admite una modificación progresivamente medible.
\end{teorema}
\subsection{Procesos de Weiner}
\begin{definicion}[Proceso de Weiner]
    \label{def:proceso_weiner}
    Un \termino[proceso!de Weiner]{proceso de Weiner} $W$ es un proceso estocástico adaptado a una filtración $\{\mathcal{F}\}_{t\in[0,T]}$ que verifica las siguientes propiedades:
    \begin{itemize}
        \item La variable aleatoria $W_0$ es nula en casi todo punto, es decir, $\Pro(W_0=0)=1$.
        \item Las trayectorias del proceso son continuas. Esto es, para cada $\omega$, su trayectoria asociada $X_\omega\colon[0,T]\to\R$ es continua.
        \item Si $0\leq s<t$, la variable aleatoria $W_t-W_s$:
              \begin{itemize}
                  \item Es independiente de la $\sigma$-álgebra $\mathcal{F}_s$.
                  \item Sigue una distribución $\Normal(0,t-s)$.
              \end{itemize}
    \end{itemize}
\end{definicion}